\documentclass[12pt]{article}
\usepackage[T1]{fontenc}
\usepackage[utf8]{inputenc}
\usepackage{lmodern}
\usepackage{textcomp}
\usepackage{enumitem}
\usepackage{geometry}
\usepackage{graphicx}
\usepackage{amsmath, amssymb}
\geometry{margin=1in}
\begin{document}
\begin{center}
Physics - Graduate Level Assessment 25 \
Total Marks: 40 \
Answer all questions.
\end{center}

\section*{Multiple Choice (10 marks)}
\begin{enumerate}[label=\arabic*.]
\item What is the refracting power of a spherical refracting surface of radius 20 cm separating air (n = 1) from glass (n = 1.5, where n denotes the index of refraction of the medium) ?
\begin{enumerate}[label=(\alph*)]
    \item 1.0 diopter
    \item 4.0 diopters
    \item 2.0 diopters
    \item 2.75 diopters
    \item 3.5 diopters
\end{enumerate}
\item The nebular theory of the formation of the solar system successfully predicts all but one of the following. Which one does the theory not predict?
\begin{enumerate}[label=(\alph*)]
    \item The presence of moons around several planets
    \item Planets orbit around the Sun in nearly circular orbits in a flattened disk.
    \item the equal number of terrestrial and jovian planets
    \item The existence of the Sun's magnetic field
    \item The different chemical compositions of the planets
\end{enumerate}
\item Under what conditions is the net electric flux through a closed surface proportional to the enclosed charge?
\begin{enumerate}[label=(\alph*)]
    \item only when there are no charges outside the surface
    \item only when the charges are in motion
    \item only when the electric field is uniform
    \item only when the enclosed charge is negative
    \item only when the charge is at rest
\end{enumerate}
\item Does an ideal gas get hotter or colder when it expands according to the law $pV^2$ = Const?
\begin{enumerate}[label=(\alph*)]
    \item remains the same
    \item none of these
    \item hotter
    \item the temperature oscillates
    \item it gets colder, but only under low pressure
\end{enumerate}
\item An object is placed 10 cm in front of a diverging mirror. What is the focal length of the mirror if the image appears 2 cm behind the mirror?
\begin{enumerate}[label=(\alph*)]
    \item -2/5 cm
    \item -3/2 cm
    \item -2/3 cm
    \item -5/2 cm
    \item -5 cm
\end{enumerate}
\end{enumerate}

\section*{True / False (10 marks)}
\textit{Answer True or False}
\begin{enumerate}[label=\arabic*.]
\setcounter{enumi}{5}
\item True or False: The thickness of the glass plate that causes the object to appear five-sixths of an inch nearer is 1.0 inch.
\item True or False: As the angle of the inclined plane increases, the normal force remains constant until a certain angle, then decreases linearly.
\item True or False: The diameter of Pluto was first accurately measured using brightness measurements during mutual eclipses of Pluto and Charon.
\item True or False: The axis of the Earth is tilted at an angle of approximately 23.5 degrees relative to its orbital plane around the Sun.
\item True or False: As a fire engine approaches, the amplitude of its siren sound increases due to the Doppler effect.
\end{enumerate}

\section*{Long Form Response (20 marks)}
\textit{Answer each question with a clear and complete explanation.}
\begin{enumerate}[label=\arabic*.]
\setcounter{enumi}{10}
\item Analyze the double-slit experiment with a slit separation of 0.1 mm and a distance to the screen of 1 m using yellow light. Calculate and discuss the implications of the spacing between adjacent bright lines.
\item Discuss the energy transition of an electron moving from the n = 3 to n = 2 orbit, specifically focusing on the quantum of energy released as visible red light. Calculate and analyze its magnitude.
\end{enumerate}

\vfill
\noindent\textit{End of Paper}
\end{document}